\subsection{Strengthening Cybersecurity in AI-Enabled Systems}
\label{sec:9.2.2}
An AI system embedded in infrastructure fails differently from ordinary software. Exploiting it need not crash anything; it can leave the system running and confidently wrong about a decision somebody trusted it to make.

The clearest demonstration is a decade old and still the one to reach for. Researchers placed a small pattern of black-and-white stickers on a real stop sign and a standard road-sign classifier read it as a speed limit sign in the large majority of frames captured from a moving vehicle \autocite{eykholt2018robust} — an instance of the vulnerability first documented in image classifiers years earlier \autocite{szegedy2013intriguing}. Nothing about the attack required access to the system. It required access to the world the system was looking at.

The defenses are the ones section~\ref{sec:6.2.2} describes, plus the ordinary security practice AI deployments often skip — adversarial training during development \autocite{goodfellow2014explaining}, anomaly detection over live behavior \autocite{darktrace2014enterprise}, red-teaming before deployment rather than after an incident, and standards written for AI systems instead of adapted afterward from general software rules. What separates them is whether the guarantee survives a stronger attacker, and that is checkable rather than a matter of taste. Defensive distillation was published as a hardening technique and its protection weakened considerably once attacks written against it were tried \autocite{carlini2016defensive}; randomized smoothing instead offers a certified bound — a radius around each input inside which the prediction provably cannot be flipped — which is a smaller promise and a kept one \autocite{cohen2019certified}.

None of it works in isolation, which is the part most often got wrong. A model hardened by adversarial training, running on infrastructure nobody monitors, is exposed; a well-monitored system built on an architecture that folds under a basic perturbation is exposed differently. Security here is a property of the whole stack, and reporting the strongest link is how a deployment gets described as secure while remaining trivially attackable.
